\chapter{Analyse des besoins et conception}
\label{ch:conception}

\section{Introduction}

Ce chapitre présente l'analyse des besoins fonctionnels de \nomapplication, l'identification des acteurs et la modélisation UML qui a guidé l'implémentation.

\section{Identification des acteurs}

Le système distingue \textbf{quatre acteurs principaux} :

\begin{table}[H]
    \centering
    \caption{Acteurs du système \nomapplication}
    \label{tab:acteurs}
    \begin{tabular}{@{}lp{9cm}@{}}
        \toprule
        \textbf{Acteur} & \textbf{Rôle et responsabilités} \\
        \midrule
        Agent bancaire & Clients, comptes, opérations courantes, historique, chéquiers \\
        Chef d'agence & Supervision, dashboard, notifications, journal d'audit \\
        Administrateur & Gestion du personnel (création, rôles, activation) \\
        Client bancaire & Portail en ligne : consultation, reçus, notifications \\
        \bottomrule
    \end{tabular}
\end{table}

L'authentification est unifiée sur \texttt{/login} ; la redirection post-connexion dépend du rôle (personnel $\rightarrow$ \texttt{/dashboard}, client $\rightarrow$ \texttt{/portal/dashboard}).

\section{Analyse des besoins fonctionnels}

\subsection{Périmètre}

\begin{itemize}[leftmargin=*]
    \item \textbf{Inclus} : CRUD clients/comptes, opérations financières, utilisateurs, audit, reporting, paiement facture, chéquier, notifications, portail client.
    \item \textbf{Exclu} : crédit, cartes bancaires, GAB, interbancaire, microservices.
\end{itemize}

\subsection{Fonctionnalités principales}

\begin{table}[H]
    \centering
    \caption{Modules fonctionnels}
    \label{tab:fonctionnalites}
    \small
    \begin{tabular}{@{}p{3cm}p{9.5cm}@{}}
        \toprule
        \textbf{Module} & \textbf{Description} \\
        \midrule
        Authentification & Login, BCrypt, rôles ADMIN / AGENT / CHEF\_AGENCE / CLIENT \\
        Clients & CRUD, recherche multicritère, statuts \\
        Comptes & Ouverture, blocage, clôture (solde nul requis) \\
        Opérations & Dépôt, retrait, virement avec règles métier \\
        Paiement facture & LYDEC, ONEE, IAM… — débit compte + reçu \\
        Chéquier & Workflow PENDING $\rightarrow$ DELIVERED, sans impact solde \\
        Reporting & Dashboard, relevé, reçu HTML + PDF (OpenPDF) \\
        Notifications & Cloche, alertes chéquier, notifications client \\
        Portail & Consultation comptes, historique, chéquiers \\
        \bottomrule
    \end{tabular}
\end{table}

\subsection{Règles de gestion}

\begin{itemize}[leftmargin=*]
    \item R1--R2 : un client possède un ou plusieurs comptes ; un compte appartient à un seul client.
    \item R3 : retrait/paiement refusé si solde insuffisant.
    \item R4 : virement — comptes source et destination actifs.
    \item R5--R8 : traçabilité, comptes bloqués/clôturés inopérants, BCrypt, audit.
    \item R9--R14 : extensions facture et chéquier (référence obligatoire, une commande PENDING max).
\end{itemize}

\section{Modélisation UML}

\subsection{Diagrammes de cas d'utilisation}

\includefigure{figures/uml/01-clients-comptes.pdf}{width=\textwidth}{01-clients-comptes}{Cas d'utilisation — clients et comptes}{fig:uc-clients}

\includefigure{figures/uml/02-operations-supervision.pdf}{width=\textwidth}{02-operations-supervision}{Cas d'utilisation — opérations et supervision}{fig:uc-ops}

\includefigure{figures/uml/03-portail-client.pdf}{width=\textwidth}{03-portail-client}{Cas d'utilisation — portail client}{fig:uc-portal}

\subsection{Diagramme de classes}

\includefigure{figures/uml/diagramme-classes.pdf}{width=0.95\textwidth}{diagramme-classes}{Diagramme de classes}{fig:classes}

Les entités principales sont : \texttt{User}, \texttt{Client}, \texttt{Account}, \texttt{Transaction}, \texttt{BillProvider}, \texttt{BillPayment}, \texttt{CheckbookOrder}, \texttt{Notification}, \texttt{AuditLog}.

\subsection{Modèle de données (MCD et MLD)}

\includefigure{figures/uml/MCD.pdf}{width=\textwidth}{MCD}{Modèle conceptuel de données (MCD)}{fig:mcd}

\includefigure{figures/uml/MLD.pdf}{width=\textwidth}{MLD}{Modèle logique de données (MLD)}{fig:mld}

\subsection{Diagrammes de séquence}

\includefigure{figures/uml/01-authentification.pdf}{width=\textwidth}{01-authentification}{Séquence — authentification dual (personnel / client)}{fig:seq-auth}

\includefigure{figures/uml/03-operations-financieres.pdf}{width=\textwidth}{03-operations-financieres}{Séquence — opérations financières}{fig:seq-ops}

\includefigure{figures/uml/06-paiement-facture.pdf}{width=\textwidth}{06-paiement-facture}{Séquence — paiement de facture}{fig:seq-facture}

\includefigure{figures/uml/07-commande-chequier.pdf}{width=\textwidth}{07-commande-chequier}{Séquence — commande de chéquier}{fig:seq-chequier}

\subsection{Diagramme d'activité}

\includefigure{figures/uml/diagramme-activite-virement.pdf}{width=0.85\textwidth}{diagramme-activite-virement}{Activité — virement bancaire}{fig:activite-virement}

\section{Conclusion}

Ce chapitre a identifié les acteurs, formalisé les besoins fonctionnels et présenté la modélisation UML complète. Le chapitre suivant détaille l'environnement technique et les outils de développement.
